\documentclass[12pt,a4paper]{article}
\usepackage[UTF8]{ctex}
\usepackage{amsmath,amssymb,amsthm}
\usepackage{geometry}
\usepackage{graphicx}
\usepackage{hyperref}
\usepackage{bm}

\geometry{margin=2.5cm}

\title{$\frac{\partial T_{0i}}{\partial \theta_j} = T_{0i} [A_j^b]$的详细证明}
\author{Modern Robotics}
\date{}

\begin{document}

\maketitle

\tableofcontents
\newpage

\section{引言}

在机器人运动学和动力学中，一个重要关系是：

\begin{equation}
\frac{\partial T_{0i}}{\partial \theta_j} = T_{0i} [A_j^b]
\label{eq:main_relation}
\end{equation}

其中：
\begin{itemize}
\item $T_{0i}(\theta_1, \ldots, \theta_i)$：从基座坐标系$\{0\}$到连杆坐标系$\{i\}$的正向运动学变换
\item $\theta_j$：关节$j$的角度（$j = 1, \ldots, i$）
\item $A_j^b$：关节$j$的螺旋轴，在坐标系$\{i\}$中表示
\item $[A_j^b]$：螺旋轴$A_j^b$的反对称矩阵表示
\end{itemize}

这个关系是理解Body雅可比的基础。本文将提供这个重要结果的详细证明。

\section{预备知识}

\subsection{正向运动学的链式结构}

正向运动学$T_{0i}$可以表示为：

\begin{equation}
T_{0i}(\theta_1, \ldots, \theta_i) = T_{01}(\theta_1) T_{12}(\theta_2) \cdots T_{i-1,i}(\theta_i)
\label{eq:forward_kinematics_chain}
\end{equation}

其中$T_{j-1,j}(\theta_j)$是关节$j$的变换矩阵。

\subsection{关节变换矩阵}

对于关节$j$，其变换矩阵为：

\begin{equation}
T_{j-1,j}(\theta_j) = M_{j-1,j} e^{[A_j]\theta_j}
\label{eq:joint_transform}
\end{equation}

其中：
\begin{itemize}
\item $M_{j-1,j}$：当$\theta_j = 0$时的初始配置
\item $A_j$：关节$j$的螺旋轴（在坐标系$\{j\}$中表示）
\item $[A_j]$：螺旋轴的反对称矩阵表示
\end{itemize}

\subsection{Twist与变换矩阵的关系}

从前面我们知道：$\dot{T} = [V]T$，其中$V$是twist。

\section{证明方法1：从链式法则出发}

\subsection{基本思路}

利用正向运动学的链式结构，对$\theta_j$求偏导时，只有包含$\theta_j$的项会变化。

\subsection{详细推导}

\textbf{步骤1}：写出正向运动学的链式结构

\begin{equation}
T_{0i} = T_{01}(\theta_1) T_{12}(\theta_2) \cdots T_{j-1,j}(\theta_j) \cdots T_{i-1,i}(\theta_i)
\end{equation}

\textbf{步骤2}：对$\theta_j$求偏导

当对$\theta_j$求偏导时，只有$T_{j-1,j}(\theta_j)$依赖于$\theta_j$，其他项都是常数：

\begin{align}
\frac{\partial T_{0i}}{\partial \theta_j} &= T_{01}(\theta_1) \cdots T_{j-2,j-1}(\theta_{j-1}) \frac{\partial T_{j-1,j}(\theta_j)}{\partial \theta_j} T_{j,j+1}(\theta_{j+1}) \cdots T_{i-1,i}(\theta_i) \\
&= T_{0,j-1} \frac{\partial T_{j-1,j}(\theta_j)}{\partial \theta_j} T_{j,i}
\label{eq:partial_chain}
\end{align}

\textbf{步骤3}：计算$\frac{\partial T_{j-1,j}(\theta_j)}{\partial \theta_j}$

从方程(\ref{eq:joint_transform})：

\begin{equation}
T_{j-1,j}(\theta_j) = M_{j-1,j} e^{[A_j]\theta_j}
\end{equation}

对$\theta_j$求导：

\begin{align}
\frac{\partial T_{j-1,j}(\theta_j)}{\partial \theta_j} &= M_{j-1,j} \frac{\partial}{\partial \theta_j}[e^{[A_j]\theta_j}] \\
&= M_{j-1,j} [A_j] e^{[A_j]\theta_j} \\
&= M_{j-1,j} [A_j] (M_{j-1,j}^{-1} T_{j-1,j}) \\
&= [A_j] T_{j-1,j}
\end{align}

\textbf{注意}：这里$[A_j]$是关节$j$的螺旋轴在坐标系$\{j\}$中的表示。

\textbf{步骤4}：转换到坐标系$\{i\}$

我们需要将$[A_j]$转换到坐标系$\{i\}$中表示。

从方程(\ref{eq:partial_chain})：

\begin{align}
\frac{\partial T_{0i}}{\partial \theta_j} &= T_{0,j-1} [A_j] T_{j-1,j} T_{j,i} \\
&= T_{0,j-1} [A_j] T_{j-1,i} \\
&= T_{0,j-1} [A_j] T_{j-1,i}
\end{align}

\textbf{关键观察}：$T_{0,j-1} [A_j] T_{j-1,i}$可以写成$T_{0i} [A_j^b]$的形式，其中$A_j^b$是$A_j$在坐标系$\{i\}$中的表示。

\textbf{步骤5}：使用伴随变换}

根据反对称矩阵的相似变换性质：

\begin{equation}
T_{0,j-1} [A_j] T_{j-1,i} = T_{0i} [\text{Ad}_{T_{i,j-1}}(A_j)]
\end{equation}

其中$\text{Ad}_{T_{i,j-1}}$是从坐标系$\{j-1\}$到$\{i\}$的伴随变换。

\textbf{定义}：$A_j^b = \text{Ad}_{T_{i,j-1}}(A_j)$，这是关节$j$的螺旋轴在坐标系$\{i\}$中的表示。

\textbf{因此}：

\begin{equation}
\frac{\partial T_{0i}}{\partial \theta_j} = T_{0i} [A_j^b]
\label{eq:result1}
\end{equation}

\textbf{证毕！}

\section{证明方法2：从$\dot{T} = [V]T$出发}

\subsection{基本思路}

利用$\dot{T} = [V]T$的关系，当只有关节$j$运动时，$T_{0i}$的变化率应该与关节$j$的螺旋轴相关。

\subsection{详细推导}

\textbf{步骤1}：考虑只有关节$j$运动的情况

假设只有关节$j$在运动，即$\dot{\theta}_j \neq 0$，而其他关节速度为零：$\dot{\theta}_k = 0$（$k \neq j$）。

\textbf{步骤2}：计算$\dot{T}_{0i}$

当只有关节$j$运动时：

\begin{equation}
\dot{T}_{0i} = \frac{\partial T_{0i}}{\partial \theta_j} \dot{\theta}_j
\end{equation}

\textbf{步骤3}：使用$\dot{T} = [V]T$的关系

根据$\dot{T} = [V]T$，当只有关节$j$运动时，连杆$i$的Body Twist为：

\begin{equation}
V_i^b = A_j^b \dot{\theta}_j
\end{equation}

其中$A_j^b$是关节$j$的螺旋轴在坐标系$\{i\}$中的表示。

\textbf{因此}：

\begin{equation}
\dot{T}_{0i} = [V_i^b] T_{0i} = [A_j^b \dot{\theta}_j] T_{0i}
\end{equation}

\textbf{步骤4}：提取$\dot{\theta}_j$

由于$[A_j^b \dot{\theta}_j] = \dot{\theta}_j [A_j^b]$（反对称矩阵的标量乘法），我们有：

\begin{equation}
\dot{T}_{0i} = \dot{\theta}_j [A_j^b] T_{0i}
\end{equation}

\textbf{步骤5}：与偏导数比较}

另一方面：

\begin{equation}
\dot{T}_{0i} = \frac{\partial T_{0i}}{\partial \theta_j} \dot{\theta}_j
\end{equation}

\textbf{因此}：

\begin{equation}
\frac{\partial T_{0i}}{\partial \theta_j} \dot{\theta}_j = \dot{\theta}_j [A_j^b] T_{0i}
\end{equation}

由于这对任意$\dot{\theta}_j$都成立，我们得到：

\begin{equation}
\frac{\partial T_{0i}}{\partial \theta_j} = [A_j^b] T_{0i}
\end{equation}

\textbf{注意}：这里需要验证顺序。实际上应该是：

\begin{equation}
\frac{\partial T_{0i}}{\partial \theta_j} = T_{0i} [A_j^b]
\end{equation}

\textbf{为什么是$T_{0i} [A_j^b]$而不是$[A_j^b] T_{0i}$？}

这取决于$A_j^b$是在哪个坐标系中表示的。如果$A_j^b$是在坐标系$\{i\}$中表示（Body形式），则：

\begin{equation}
T_{0i}^{-1} \dot{T}_{0i} = [A_j^b \dot{\theta}_j]
\end{equation}

因此：

\begin{equation}
\dot{T}_{0i} = T_{0i} [A_j^b \dot{\theta}_j] = T_{0i} [A_j^b] \dot{\theta}_j
\end{equation}

所以：

\begin{equation}
\frac{\partial T_{0i}}{\partial \theta_j} = T_{0i} [A_j^b]
\end{equation}

\textbf{证毕！}

\section{证明方法3：从指数映射出发}

\subsection{基本思路}

利用$SE(3)$的指数映射性质，直接计算偏导数。

\subsection{详细推导}

\textbf{步骤1}：正向运动学的指数形式

正向运动学可以写成：

\begin{equation}
T_{0i}(\theta) = \prod_{k=1}^{i} M_{k-1,k} e^{[A_k]\theta_k}
\end{equation}

\textbf{步骤2}：对$\theta_j$求偏导}

当对$\theta_j$求偏导时，只有包含$e^{[A_j]\theta_j}$的项会变化：

\begin{align}
\frac{\partial T_{0i}}{\partial \theta_j} &= \left(\prod_{k=1}^{j-1} M_{k-1,k} e^{[A_k]\theta_k}\right) \frac{\partial}{\partial \theta_j}[M_{j-1,j} e^{[A_j]\theta_j}] \left(\prod_{k=j+1}^{i} M_{k-1,k} e^{[A_k]\theta_k}\right) \\
&= T_{0,j-1} \frac{\partial}{\partial \theta_j}[M_{j-1,j} e^{[A_j]\theta_j}] T_{j,i}
\end{align}

\textbf{步骤3}：计算$\frac{\partial}{\partial \theta_j}[M_{j-1,j} e^{[A_j]\theta_j}]$

\begin{align}
\frac{\partial}{\partial \theta_j}[M_{j-1,j} e^{[A_j]\theta_j}] &= M_{j-1,j} [A_j] e^{[A_j]\theta_j} \\
&= [A_j] M_{j-1,j} e^{[A_j]\theta_j} \\
&= [A_j] T_{j-1,j}
\end{align}

\textbf{步骤4}：组合}

\begin{align}
\frac{\partial T_{0i}}{\partial \theta_j} &= T_{0,j-1} [A_j] T_{j-1,j} T_{j,i} \\
&= T_{0,j-1} [A_j] T_{j-1,i}
\end{align}

\textbf{步骤5}：转换到Body形式}

使用反对称矩阵的相似变换：

\begin{equation}
T_{0,j-1} [A_j] T_{j-1,i} = T_{0i} T_{i,j-1} [A_j] T_{j-1,i} = T_{0i} [\text{Ad}_{T_{i,j-1}}(A_j)]
\end{equation}

定义$A_j^b = \text{Ad}_{T_{i,j-1}}(A_j)$，则：

\begin{equation}
\frac{\partial T_{0i}}{\partial \theta_j} = T_{0i} [A_j^b]
\end{equation}

\textbf{证毕！}

\section{物理意义}

\subsection{直观理解}

\textbf{为什么是$T_{0i} [A_j^b]$？}

\begin{itemize}
\item $A_j^b$是关节$j$的螺旋轴在坐标系$\{i\}$中的表示（Body形式）
\item 当只有关节$j$运动时，它产生的运动在坐标系$\{i\}$中观察就是$A_j^b \dot{\theta}_j$
\item 根据$\dot{T} = [V]T$，如果twist是$A_j^b \dot{\theta}_j$（在坐标系$\{i\}$中表示），则：
\begin{equation}
T_{0i}^{-1} \dot{T}_{0i} = [A_j^b \dot{\theta}_j]
\end{equation}
\item 因此：$\dot{T}_{0i} = T_{0i} [A_j^b \dot{\theta}_j] = T_{0i} [A_j^b] \dot{\theta}_j$
\item 所以：$\frac{\partial T_{0i}}{\partial \theta_j} = T_{0i} [A_j^b]$
\end{itemize}

\subsection{Body形式与Spatial形式}

\textbf{Body形式}（在坐标系$\{i\}$中表示）：
\begin{equation}
\frac{\partial T_{0i}}{\partial \theta_j} = T_{0i} [A_j^b]
\end{equation}

\textbf{Spatial形式}（在固定坐标系$\{0\}$中表示）：
\begin{equation}
\frac{\partial T_{0i}}{\partial \theta_j} = [A_j^s] T_{0i}
\end{equation}

其中$A_j^s$是关节$j$的螺旋轴在固定坐标系中的表示。

\textbf{关系}：$A_j^s = \text{Ad}_{T_{0i}}(A_j^b)$

\section{具体例子}

\subsection{例子1：单关节情况}

考虑$T_{01}(\theta_1)$，即从基座到第一个连杆的变换。

\textbf{关节变换}：
\begin{equation}
T_{01}(\theta_1) = M_{01} e^{[A_1]\theta_1}
\end{equation}

\textbf{对$\theta_1$求导}：
\begin{align}
\frac{\partial T_{01}}{\partial \theta_1} &= M_{01} [A_1] e^{[A_1]\theta_1} \\
&= [A_1] M_{01} e^{[A_1]\theta_1} \\
&= [A_1] T_{01}
\end{align}

\textbf{注意}：这里$A_1$在坐标系$\{1\}$中表示，所以这是Body形式。

如果写成$T_{01} [A_1^b]$的形式，其中$A_1^b = A_1$（因为已经在坐标系$\{1\}$中），则：

\begin{equation}
\frac{\partial T_{01}}{\partial \theta_1} = T_{01} [A_1^b]
\end{equation}

\textbf{验证}：$T_{01} [A_1^b] = T_{01} [A_1] = [A_1] T_{01}$（因为$[A_1]$在Body坐标系中表示）。

\subsection{例子2：两关节情况}

考虑$T_{02}(\theta_1, \theta_2) = T_{01}(\theta_1) T_{12}(\theta_2)$。

\textbf{对$\theta_1$求偏导}：

\begin{align}
\frac{\partial T_{02}}{\partial \theta_1} &= \frac{\partial T_{01}}{\partial \theta_1} T_{12} \\
&= T_{01} [A_1] T_{12} \\
&= T_{02} T_{12}^{-1} [A_1] T_{12} \\
&= T_{02} [\text{Ad}_{T_{21}}(A_1)] \\
&= T_{02} [A_1^b]
\end{align}

其中$A_1^b = \text{Ad}_{T_{21}}(A_1)$是关节1的螺旋轴在坐标系$\{2\}$中的表示。

\textbf{对$\theta_2$求偏导}：

\begin{align}
\frac{\partial T_{02}}{\partial \theta_2} &= T_{01} \frac{\partial T_{12}}{\partial \theta_2} \\
&= T_{01} T_{12} [A_2] \\
&= T_{02} [A_2^b]
\end{align}

其中$A_2^b = A_2$（因为$A_2$已经在坐标系$\{2\}$中表示）。

\section{在Body雅可比构造中的应用}

\subsection{Body雅可比的构造}

从$\frac{\partial T_{0i}}{\partial \theta_j} = T_{0i} [A_j^b]$，我们可以得到：

\begin{equation}
T_{0i}^{-1} \frac{\partial T_{0i}}{\partial \theta_j} = [A_j^b]
\end{equation}

\textbf{因此}：Body雅可比$J_i^b(\theta)$的第$j$列就是$A_j^b$（$j = 1, \ldots, i$）。

\subsection{完整Body雅可比}

\begin{equation}
J_i^b(\theta) = \begin{bmatrix} A_1^b & A_2^b & \cdots & A_i^b & 0 & \cdots & 0 \end{bmatrix}
\end{equation}

其中：
\begin{itemize}
\item 前$i$列：各个关节的螺旋轴在坐标系$\{i\}$中的表示
\item 后$n-i$列：全为零（因为后$n-i$个关节不影响连杆$i$）
\end{itemize}

\section{总结}

\subsection{关键结论}

\begin{enumerate}
\item \textbf{基本关系}：$\frac{\partial T_{0i}}{\partial \theta_j} = T_{0i} [A_j^b]$

\item \textbf{物理意义}：
\begin{itemize}
\item 当只有关节$j$运动时，$T_{0i}$的变化率由关节$j$的螺旋轴决定
\item $A_j^b$是关节$j$的螺旋轴在坐标系$\{i\}$中的表示（Body形式）
\item 这个关系反映了关节运动对正向运动学的影响
\end{itemize}

\item \textbf{在Body雅可比中的应用}：
\begin{itemize}
\item $T_{0i}^{-1} \frac{\partial T_{0i}}{\partial \theta_j} = [A_j^b]$
\item Body雅可比的第$j$列就是$A_j^b$
\item 这提供了计算Body雅可比的方法
\end{itemize}
\end{enumerate}

\subsection{证明方法总结}

\begin{itemize}
\item \textbf{方法1}：从链式法则出发，逐步计算
\item \textbf{方法2}：从$\dot{T} = [V]T$的关系出发
\item \textbf{方法3}：从指数映射出发
\end{itemize}

所有方法都证明了相同的结果。

\section{参考文献}

\begin{itemize}
\item Modern Robotics: Mechanics, Planning, and Control
\item 第3章：刚体运动（Twist和Screw理论）
\item 第8章：开链动力学
\end{itemize}

\end{document}

